%2multibyte Version: 5.50.0.2960 CodePage: 1252

\documentclass[12pt]{article}
%%%%%%%%%%%%%%%%%%%%%%%%%%%%%%%%%%%%%%%%%%%%%%%%%%%%%%%%%%%%%%%%%%%%%%%%%%%%%%%%%%%%%%%%%%%%%%%%%%%%%%%%%%%%%%%%%%%%%%%%%%%%%%%%%%%%%%%%%%%%%%%%%%%%%%%%%%%%%%%%%%%%%%%%%%%%%%%%%%%%%%%%%%%%%%%%%%%%%%%%%%%%%%%%%%%%%%%%%%%%%%%%%%%%%%%%%%%%%%%%%%%%%%%%%%%%
\usepackage{amsfonts}
\usepackage{amsmath}

\setcounter{MaxMatrixCols}{10}
%TCIDATA{OutputFilter=LATEX.DLL}
%TCIDATA{Version=5.50.0.2960}
%TCIDATA{Codepage=1252}
%TCIDATA{<META NAME="SaveForMode" CONTENT="1">}
%TCIDATA{BibliographyScheme=Manual}
%TCIDATA{Created=Monday, September 22, 2003 09:40:11}
%TCIDATA{LastRevised=Friday, October 31, 2025 13:08:20}
%TCIDATA{<META NAME="GraphicsSave" CONTENT="32">}
%TCIDATA{<META NAME="DocumentShell" CONTENT="General\SW\Blank - Standard LaTeX Article">}
%TCIDATA{Language=American English}
%TCIDATA{CSTFile=LaTeX article (bright).cst}

\newtheorem{theorem}{Theorem}
\newtheorem{acknowledgement}[theorem]{Acknowledgement}
\newtheorem{algorithm}[theorem]{Algorithm}
\newtheorem{axiom}[theorem]{Axiom}
\newtheorem{case}[theorem]{Case}
\newtheorem{claim}[theorem]{Claim}
\newtheorem{conclusion}[theorem]{Conclusion}
\newtheorem{condition}[theorem]{Condition}
\newtheorem{conjecture}[theorem]{Conjecture}
\newtheorem{corollary}[theorem]{Corollary}
\newtheorem{criterion}[theorem]{Criterion}
\newtheorem{definition}[theorem]{Definition}
\newtheorem{example}[theorem]{Example}
\newtheorem{exercise}[theorem]{Exercise}
\newtheorem{lemma}[theorem]{Lemma}
\newtheorem{notation}[theorem]{Notation}
\newtheorem{problem}[theorem]{Problem}
\newtheorem{proposition}[theorem]{Proposition}
\newtheorem{remark}[theorem]{Remark}
\newtheorem{solution}[theorem]{Solution}
\newtheorem{summary}[theorem]{Summary}
\newenvironment{proof}[1][Proof]{\textbf{#1.} }{\ \rule{0.5em}{0.5em}}
\input{tcilatex}
\input{tcilatex}
\setlength{\textwidth}{6.5in}
\setlength{\textheight}{9in}
\setlength{\topmargin}{-.50in}
\setlength{\oddsidemargin}{.0in}
\setlength{\evensidemargin}{.0in}

\begin{document}


\begin{center}
{\LARGE Assignment 5}
\end{center}

\noindent Blankenau\hspace*{\stretch{0.7500}}

\noindent Economics 905, Fall 2025

\noindent \textit{Instructions.} You may work in groups and may compare
answers prior to submission. Joint submissions are allowed and encouraged.
All work should be legible and concise. A due date will be announced in
class.\medskip

\begin{enumerate}
\item Consider an economy consisting of overlapping generations of
two-period-lived homogeneous agents, a representative firm producing a
single good and a government. A representative agent is born each period and
the population of each generation is normalized to 1. The initial young and
old agents are endowed with human capital $h_{0}.$ All agents have one unit
of time to supply exogenously to the labor market in each period of life.
Agents beyond the initial period receive an endowment of human capital which
depends on government and family spending on education as the first period
of life begins. Specifically, the human capital of the representative agent
born in period $t$ is $h_{t}=f_{t}^{\alpha }g_{t}^{1-\alpha }$ where $f_{t}$
is spending by the agent's parents on education in period $t$ and $g_{t}$ is
spending by the government\ on education in period $t$. Assume $0\leq \alpha
\leq 1.$ In the second period of life, through experience, the agent's level
of human capital grows to $\gamma h_{t}$ where $\gamma >1.$ Parents value
consumption in each period of life. They also value the human capital of
their children. Preferences are given by 
\begin{equation*}
U_{t}=\ln c_{t,t}+\beta \ln c_{t,t+1}+\psi \ln h_{t+1}
\end{equation*}%
where $U_{t}$ is the lifetime utility of an agent born in period $t$, $%
c_{t,t}$ and $c_{t,t+1}$ are this agent's consumption when young and old and 
$h_{t+1}$ is the human capital of the agent's children. There is no storage
or capital in this economy so the structure precludes borrowing and lending.
Output is produced according to $Y_{t}=AH_{t}$ where $H_{t}$ is the amount
of human capital employed by the firm and $A>0$ is a scalar. Each period $t$
agent inelastically provides $h_{t}$ units of human capital per unit of time
when young and $\gamma h_{t}$ when old so that income is $\omega _{t}h_{t}$
and $\omega _{t+1}\gamma h_{t}$ when young and old where $\omega _{t}$ is
the wage per unit of human capital in period $t$. Government taxes
consumption while young to  pay for its education expenditure subject to a
balanced budget. Thus the first period budget constraint is $c_{t,t}\left(
1+\tau \right) =\omega _{t}h_{t}.$

\begin{enumerate}
\item Write down the problem for a period $t$ agent and find the first order
condition. The agent takes $h_{t}$, $\omega _{t},$ $\omega _{t+1}$ and the
tax rate $\tau $ as given. Also write down the firm's problem and give the
government's budget constraint.

\item Solve for the steady state level of $f$ in this economy. How does
government education spending affect total education spending?

\item Now suppose the situation is the same except that young agents can
save. Let $s_{t}$ be the representative agent's choice of savings while
young and $Ds_{t}$ $(D>0)$ \ be the income generated from this while old .
This would arise if $Y_{t}=AH_{t}+DK_{t}$ where $K_{t}$ is capital and $%
s_{t-1}=K_{t}$ in equilibrium. However, those particulars will not be
important for your answer. Let $\gamma =\beta =1$ to simplify. Find the
budget constraints while young and old assuming again that consumption while
young is taxed at rate $\tau \ $and find first order condition(s).

\item Solve for the steady state level of $f$ in this economy. Compare $f$
in this case with your finding in the earlier case and give the intuition
for any differences or why they are the same.
\end{enumerate}
\end{enumerate}

\pagebreak

\begin{enumerate}
\item 
\begin{enumerate}
\item ...%
\begin{equation*}
\underset{f_{t+1}}{\max }\ln \left( \omega _{t}h_{t}-\tau _{t,t}\right)
+\beta \ln \left( \omega _{t+1}h_{t}\gamma -f_{t+1}\right) +\psi \ln \left(
f_{t+1}^{\alpha }g_{t+1}^{1-\alpha }\right) 
\end{equation*}%
\begin{equation*}
g_{t+1}=\tau _{t,t}
\end{equation*}%
firm%
\begin{equation*}
\underset{H_{t}}{\max }AH_{t}-\omega _{t}H_{t}
\end{equation*}

\item 
\begin{equation*}
\frac{\beta }{\left( \omega _{t+1}h_{t}\gamma -f_{t+1}\right) }+\alpha \psi 
\frac{f_{t+1}^{\alpha -1}g_{t+1}^{1-\alpha }}{f_{t+1}^{\alpha
}g_{t+1}^{1-\alpha }}
\end{equation*}%
In a steady state%
\begin{equation*}
\frac{\beta }{Af^{\alpha }g^{1-\alpha }\gamma -f}=\frac{\alpha \psi }{f}
\end{equation*}%
from which 
\begin{equation*}
\beta f=\alpha \psi \left( Af^{\alpha }g^{1-\alpha }\gamma -f\right) 
\end{equation*}%
and solving gives%
\begin{equation*}
\left( \frac{A\alpha \psi \gamma }{\beta +\alpha \psi }\right) ^{\frac{1}{%
1-\alpha }}g=f.
\end{equation*}

\item With $\gamma =\beta =1,$ now the problem is 
\begin{equation*}
\underset{f_{t+1},s_{t}}{\max }\ln \left( \frac{\omega _{t}h_{t}-s_{t}}{%
1+\tau }\right) +\ln \left( \omega _{t+1}h_{t}-f_{t+1}+Ds_{t}\right) +\psi
\left( f_{t+1}^{\alpha }g_{t+1}^{1-\alpha }\right) 
\end{equation*}%
and first order conditions are%
\begin{equation*}
\frac{1}{\omega _{t}h_{t}-s_{t}}=\frac{D}{\omega _{t+1}h_{t}-f_{t+1}+Ds_{t}}
\end{equation*}%
\begin{equation*}
\frac{1}{\omega _{t+1}h_{t}-f_{t+1}+Ds_{t}}=\frac{\alpha \psi }{f_{t+1}}
\end{equation*}

\item In a steady state%
\begin{equation*}
\frac{1}{Ah-s}=\frac{D}{Ah-f+Ds}
\end{equation*}%
$\allowbreak $ $\allowbreak $%
\begin{equation*}
\frac{1}{Ah-f+Ds}=\frac{\alpha \psi }{f}
\end{equation*}%
$\allowbreak $From the first%
\begin{equation*}
s=\frac{1}{2D}\left( f-Ah+AhD\right) 
\end{equation*}%
This into the second gives%
\begin{equation*}
\frac{1}{Ah-f+Ds}=\frac{\alpha \psi }{f}
\end{equation*}%
\begin{equation*}
\frac{1}{Ah-f+\frac{1}{2}\left( f-Ah+AhD\right) }=\frac{\alpha \psi }{f}
\end{equation*}%
, Solution is: $\left\{ 
\begin{array}{ccc}
\emptyset  & \text{if} & \alpha \psi +2=0\wedge Ah\alpha \psi +AhD\alpha
\psi \neq 0\vee \frac{1}{\alpha \psi +2}\left( Ah\alpha \psi +AhD\alpha \psi
\right) \in \left\{ 0,Ah\left( D+1\right) \right\} \wedge \alpha \psi +2\neq
0 \\ 
\left\{ \frac{1}{\alpha \psi +2}\left( Ah\alpha \psi +AhD\alpha \psi \right)
\right\}  & \text{if} & \frac{1}{\alpha \psi +2}\left( Ah\alpha \psi
+AhD\alpha \psi \right) \in 
%TCIMACRO{\U{2102} }%
%BeginExpansion
\mathbb{C}
%EndExpansion
\setminus \left\{ 0,Ah\left( D+1\right) \right\} \wedge \alpha \psi +2\neq 0
\\ 
%TCIMACRO{\U{2102} }%
%BeginExpansion
\mathbb{C}
%EndExpansion
\setminus \left\{ 0,Ah+AhD\right\}  & \text{if} & \alpha \psi +2=0\wedge
Ah\alpha \psi +AhD\alpha \psi =0%
\end{array}%
\right. \allowbreak \allowbreak $%
\begin{equation*}
f=h\frac{\left( Ah\alpha \psi +AhD\alpha \psi \right) }{\alpha \psi +2}
\end{equation*}%
\begin{equation*}
f=\frac{\alpha \psi A\left( 1+D\right) f^{\alpha }g^{1-\alpha }}{\alpha \psi
+2}
\end{equation*}%
\begin{equation*}
f=\left( \frac{\alpha \psi A\left( 1+D\right) }{\alpha \psi +2}\right) ^{%
\frac{1}{1-\alpha }}g
\end{equation*}%
Clearly increases with $D.$
\end{enumerate}
\end{enumerate}

\end{document}
